\begin{table}[H]
  \begin{center}
    \setlength\extrarowheight{-6pt}
    \begin{tabular}{c|ccccccccc}
      $n/k$ & 0 & 1 & 2 & 3 & 4 & 5 & 6 & 7 & 8 \\ [3px]
      \hline
      0 & 1 &   &   &   &   &   &   &   &   \\
      1 & 1 & 1 &   &   &   &   &   &   &   \\
      2 & 1 & 2 & 1 &   &   &   &   &   &   \\
      3 & 1 & $\tfrac{13}{4}$ & 3 & 1 &   &   &   &   &   \\
      4 & 1 & 5 & 7 & 4 & 1 &   &   &   &   \\
      5 & 1 & $\tfrac{121}{16}$ & 15 & $\tfrac{25}{2}$ & 5 & 1 &   &   &   \\
      6 & 1 & $\tfrac{91}{8}$ & 31 & 35 & 20 & 6 & 1 &   &   \\
      7 & 1 & $\tfrac{1093}{64}$ & 63 & $\tfrac{1491}{16}$ & 70 & $\tfrac{119}{4}$ & 7 & 1 &   \\
      8 & 1 & $\tfrac{205}{8}$ & 127 & $\tfrac{483}{2}$ & 231 & 126 & 42 & 8 & 1 \\
    \end{tabular}
  \end{center}
  \caption{
    Values of generalized central factorial numbers $T_1(n,k)$.
    See also the sequences \href{https://oeis.org/A395860}{\texttt{A395860}},
    and \href{https://oeis.org/A395861}{\texttt{A395861}} in the OEIS~\cite{sloane2003line}.
  }
  \label{tab:central-factorial-numbers-t1}
\end{table}
